\documentclass[a4paper,11pt]{article}

\usepackage[margin=2.5cm]{geometry}
\usepackage[T1]{fontenc}
\usepackage[utf8]{inputenc}
\usepackage[english]{babel}
\usepackage{lmodern}
\usepackage{microtype}
\usepackage{amsmath}
\usepackage{booktabs}
\usepackage{tabularx}
\usepackage{float}
\usepackage{hyperref}
\usepackage{xcolor}

\hypersetup{
  colorlinks=true,
  linkcolor=blue!50!black,
  urlcolor=blue!50!black,
  pdftitle={IoT Challenge 2 -- Part 2 (Exercise)},
  pdfauthor={Tommaso Neri}
}

\setlength{\parindent}{0pt}
\setlength{\parskip}{0.55em}
\renewcommand{\arraystretch}{1.18}
\linespread{1.04}

\newcommand{\coursename}{Internet of Things Lab}
\newcommand{\challengename}{IoT Challenge \#2 -- Part 2 (Exercise)}
\newcommand{\authorname}{Tommaso Neri}
\newcommand{\personcode}{10817232}
\newcommand{\answer}[1]{\textbf{\textcolor{red}{#1}}}
\newcommand{\tx}{\mathrm{TX}}
\newcommand{\rx}{\mathrm{RX}}
\newcommand{\uJ}{\mu\mathrm{J}}

\title{\challengename}
\author{
  \authorname\\
  Person Code: \personcode
}
\date{\today}

\begin{document}

\maketitle
\thispagestyle{empty}

\section*{General Parameters \& Constants}
The sensor transmits one measurement every 2 minutes. Therefore, in the
1-hour observation window:
\[
  N = \frac{60\ \mathrm{min}}{2\ \mathrm{min/update}} = 30
  \quad \text{updates/hour}.
\]
I assume the first transmission occurs at $t=2$ minutes, so the 1-hour window
contains 30 generated measurements.

All message sizes in the exercise are given in bytes, while the radio energy
costs are given per bit. Hence, for a message of length $L$ bytes:
\[
  E_{\tx}(L) = 8L \cdot 50 = 400L\ \mathrm{nJ},
  \qquad
  E_{\rx}(L) = 8L \cdot 58 = 464L\ \mathrm{nJ}.
\]

The costs per byte are therefore:
\begin{itemize}
    \item $E_{\tx} = 50\ \mathrm{nJ/bit} = 400\ \mathrm{nJ/byte}$
    \item $E_{\rx} = 58\ \mathrm{nJ/bit} = 464\ \mathrm{nJ/byte}$
\end{itemize}

Only the battery-powered sensor and actuator are considered; the gateway is
mains-powered and its energy consumption is ignored. Since the communication is
ideal, no retransmission term is needed.

\section*{EQ1: Energy Consumption over 1 Hour}

\subsection*{(a) Case 1: CoAP (Direct Communication)}
CoAP is used directly between the sensor and the actuator. For each update, the
sensor transmits one Confirmable PUT Request of 70 bytes and receives one PUT
ACK of 15 bytes.

The energy consumed by the sensor over the 30 updates is:
\[
\begin{aligned}
E_{\mathrm{Sensor,CoAP}}
  &= N \cdot \left(E_{\tx}(70) + E_{\rx}(15)\right) \\
  &= 30 \cdot (70 \cdot 400\ \mathrm{nJ} + 15 \cdot 464\ \mathrm{nJ}) \\
  &= 30 \cdot (28{,}000\ \mathrm{nJ} + 6{,}960\ \mathrm{nJ}) \\
  &= 30 \cdot 34{,}960\ \mathrm{nJ} \\
  &= 1{,}048{,}800\ \mathrm{nJ} \\
  &= 1048.8\ \uJ.
\end{aligned}
\]

The actuator (subscriber) receives the CON PUT Request and transmits a PUT ACK directly to the sensor.
The energy consumed by the actuator is:
\[
\begin{aligned}
E_{\mathrm{Actuator,CoAP}}
  &= N \cdot \left(E_{\rx}(70) + E_{\tx}(15)\right) \\
  &= 30 \cdot (70 \cdot 464\ \mathrm{nJ} + 15 \cdot 400\ \mathrm{nJ}) \\
  &= 30 \cdot (32{,}480\ \mathrm{nJ} + 6{,}000\ \mathrm{nJ}) \\
  &= 30 \cdot 38{,}480\ \mathrm{nJ} \\
  &= 1{,}154{,}400\ \mathrm{nJ} \\
  &= 1154.4\ \uJ.
\end{aligned}
\]

Overall, the energy consumed using CoAP is the sum of the two values:
\[
  E_{\mathrm{CoAP}}
  = E_{\mathrm{Sensor,CoAP}} + E_{\mathrm{Actuator,CoAP}}
  = 2203.2\ \uJ.
\]

\textbf{Answer EQ1(a):} \answer{2203.2 $\mu\mathrm{J}$}

\subsection*{(b) Case 2: MQTT (Via Gateway)}
For EQ1, I treat the 1-hour observation window as one active MQTT session for
both battery-powered devices. Devices start disconnected, so each MQTT client
connects once at the beginning and then stays connected for the whole hour. No
DISCONNECT message is counted in EQ1, and the actuator subscribes only once. The
explicit sleep/wake behavior is evaluated separately in EQ2. Whether the MQTT
session is persistent or clean does not change this EQ1 energy count.

The sensor first performs CONNECT/CONNACK. Then, for each update, it sends a
QoS 1 PUBLISH to the gateway and receives the PUBACK:
\[
\begin{aligned}
E_{\mathrm{S,conn}}
  &= E_{\tx}(60) + E_{\rx}(50) \\
  &= 60 \cdot 400 + 50 \cdot 464 \\
  &= 47{,}200\ \mathrm{nJ}.
\end{aligned}
\]

\[
\begin{aligned}
E_{\mathrm{S,pub}}
  &= E_{\tx}(75) + E_{\rx}(50) \\
  &= 75 \cdot 400 + 50 \cdot 464 \\
  &= 53{,}200\ \mathrm{nJ/update}.
\end{aligned}
\]

\[
\begin{aligned}
E_{\mathrm{Sensor,MQTT}}
  &= E_{\mathrm{S,conn}} + N \cdot E_{\mathrm{S,pub}} \\
  &= 47{,}200 + 30 \cdot 53{,}200 \\
  &= 1{,}643{,}200\ \mathrm{nJ}
   = 1643.2\ \uJ.
\end{aligned}
\]

The actuator first performs CONNECT/CONNACK and SUBSCRIBE/SUBACK. Then, for
each update, it receives a QoS 1 PUBLISH from the gateway and transmits the
PUBACK:
\[
\begin{aligned}
E_{\mathrm{A,conn}}
  &= E_{\tx}(60) + E_{\rx}(50) \\
  &= 60 \cdot 400 + 50 \cdot 464 \\
  &= 47{,}200\ \mathrm{nJ}.
\end{aligned}
\]

\[
\begin{aligned}
E_{\mathrm{A,sub}}
  &= E_{\tx}(65) + E_{\rx}(55) \\
  &= 65 \cdot 400 + 55 \cdot 464 \\
  &= 51{,}520\ \mathrm{nJ}.
\end{aligned}
\]

\[
\begin{aligned}
E_{\mathrm{A,recv}}
  &= E_{\rx}(75) + E_{\tx}(50) \\
  &= 75 \cdot 464 + 50 \cdot 400 \\
  &= 54{,}800\ \mathrm{nJ/update}.
\end{aligned}
\]

\[
\begin{aligned}
E_{\mathrm{Actuator,MQTT}}
  &= E_{\mathrm{A,conn}} + E_{\mathrm{A,sub}} + N \cdot E_{\mathrm{A,recv}} \\
  &= 47{,}200 + 51{,}520 + 30 \cdot 54{,}800 \\
  &= 1{,}742{,}720\ \mathrm{nJ}
   = 1742.72\ \uJ.
\end{aligned}
\]

Overall, the energy consumed using MQTT is:
\[
  E_{\mathrm{MQTT}}
  = E_{\mathrm{Sensor,MQTT}} + E_{\mathrm{Actuator,MQTT}}
  = 3385.92\ \uJ.
\]

\textbf{Answer EQ1(b):} \answer{3385.92 $\mu\mathrm{J}$}

\section*{EQ2: Actuator wakes up only once every 30 minutes}

\textbf{Which protocol would you choose for long-term operation?} \\
\answer{MQTT}

\textbf{Answer and justification for form/report (290 characters):} \\
MQTT QoS 1, sensor PUBLISH RETAIN=1, actuator Clean Session=True. Optimize actuator energy and freshness: broker keeps only the latest value; at each wake-up actuator connects, subscribes, receives one retained PUBLISH and ACKs. Clean session avoids queued old QoS1 messages while sleeping.

\end{document}
